\documentclass[11pt]{amsart}

\usepackage{amsmath}
\usepackage{amssymb}
\usepackage{amsthm}
%\usepackage{psfig}

\begin{document}
\baselineskip=16pt
\textheight=8.5in
%\parindent=0pt 
\def\sk {\hskip .5cm}
\def\skv {\vskip .08cm}
\def\cos {\mbox{cos}}
\def\sin {\mbox{sin}}
\def\tan {\mbox{tan}}
\def\intl{\int\limits}
\def\lm{\lim\limits}
\newcommand{\frc}{\displaystyle\frac}
\def\xbf{{\mathbf x}}
\def\fbf{{\mathbf f}}
\def\gbf{{\mathbf g}}

\def\dbA{{\mathbb A}}
\def\dbB{{\mathbb B}}
\def\dbC{{\mathbb C}}
\def\dbD{{\mathbb D}}
\def\dbE{{\mathbb E}}
\def\dbF{{\mathbb F}}
\def\dbG{{\mathbb G}}
\def\dbH{{\mathbb H}}
\def\dbI{{\mathbb I}}
\def\dbJ{{\mathbb J}}
\def\dbK{{\mathbb K}}
\def\dbL{{\mathbb L}}
\def\dbM{{\mathbb M}}
\def\dbN{{\mathbb N}}
\def\dbO{{\mathbb O}}
\def\dbP{{\mathbb P}}
\def\dbQ{{\mathbb Q}}
\def\dbR{{\mathbb R}}
\def\dbS{{\mathbb S}}
\def\dbT{{\mathbb T}}
\def\dbU{{\mathbb U}}
\def\dbV{{\mathbb V}}
\def\dbW{{\mathbb W}}
\def\dbX{{\mathbb X}}
\def\dbY{{\mathbb Y}}
\def\dbZ{{\mathbb Z}}

\def\la{{\langle}}
\def\ra{{\rangle}}
\def\summ{{\sum\limits}}
\def\eps{{\varepsilon}}
\def\lam{{\lambda}}
\def\tr{{\rm tr}}
\def\char{{\rm char}}
\def\Mat{{\rm Mat}}
\def\Func{{\rm Func}}
\def\Span{{\rm Span}}



\bf\centerline{Homework \#2. Due on Thursday, Sep 10th, 11:59pm on Canvas}\rm
\vskip .1cm

\bf{Reading for this homework: }\rm Sections 1.4-1.6

\bf{Plan for next week (Sep 7,9): }\rm  1.6 (bases and dimension), 2.1 (linear transformations, kernels and images), maybe start 2.2
(matrix of a linear transformation). 
\skv
In this homework you can use Theorems~A and B below. They are both proved in 1.6 (Theorem~B is stated slightly differently) and will also be proved in class on Mon, Sep 7.

Let $V$ be a vector space over a field $F$. In class we defined $V$ to be {\it finite-dimensional} if $V$ has a finite spanning set.
We then immediately proved that if $V$ is finite-dimensional, then $V$ has a finite basis (Lemma~3.5 from class).
\skv
{\bf Theorem A: }\rm Let $V$ be a finite-dimensional vector space. Then any two bases of $V$ have the same (finite) cardinality. This common cardinality is called the dimension of $V$.
\skv
{\bf Theorem B: }\rm Let $V$ be a vector space of finite dimension $n$, and let $S$ be a subset of $V$ with $|S|=n$. Then
$S$ is a basis $\iff$ $S$ is spanning $\iff$ $S$ is linearly independent. 
\skv
\bf\centerline{Problems: }\rm
\skv

{\bf Problem 1:} Let $V$ be a vector space over a field $F$.
\begin{itemize}
\item[(a)] Prove Lemma 3.4(1) from class: If $S$ is a linearly independent subset of $V$ which is not a basis, there exists a
linearly independent subset $T$ of $V$ which strictly contains $S$ (that is, $T\supseteq S$ and $T\neq S$). {\bf Hint:} This
can be proved similarly to one of the directions in Lemma 3.2 from class.
\item[(b)] Assume that $V$ is infinite-dimensional (that is, $V$ is not finite-dimensional, as defined at the beginning of this homework).
Use (a) to prove that $V$ contains an infinite linearly independent subset. 
\end{itemize}

\skv
Before the next problem we introduce the notion of characteristic of a field. Given a field $F$ and $n\in\dbN$, we define $n_F\in F$
by $n_F=\underbrace{1_F+\ldots+1_F}_{n\mbox{ times }}$ ($1_F$ added to itself $n$ times). If $n_F\neq 0_F$ for all $n\in\dbN$, one
defines the characteristic of $F$ (denoted $\char(F)$) to be $0$; otherwise, $\char(F)$ is the smallest $n$ such that $n_F=0_F$.
For example, $\char(\dbQ)=\char(\dbR)=\char(\dbC)=0$ while $\char(\dbZ_p)=p$ for any prime $p$. It is not difficult to show that
$\char(F)$ is either $0$ or a prime number for any field $F$.

 \skv 
{\bf Problem 2:} Let $F$ be a field, $V$ a vector space over $F$ and $u,v\in V$ distinct elements such that $\{u,v\}$ is linearly independent. 
\skv
\begin{itemize}
\item[(a)] Prove that $u+v\neq v$ and
$\{u+v,v\}$ is linearly independent.
\item[(b)] Assume that $\char(F)\neq 2$. Prove that $u+v\neq u-v$ and
$\{u+v,u-v\}$ is linearly independent. 

Note that if $\char(F)=2$, then $v+v=0_V$ for any $v$ (since $v+v=1_F\cdot v +1_F\cdot v=
(1_F+1_F)\cdot v=2_F\cdot v=0_F\cdot v=0_V$), so $-v=v$ and hence $u-v=u+v$. 
\end{itemize}


\skv 
{\bf Problem 3:} Let $F$ be a field, $n\in\dbN$ and $\mathcal P_n(F)$ the vector space of all polynomials of degree $\leq n$
with coefficients in $F$. Let $S$ be a subset of $\mathcal P_n(F)$ such that $0\not\in S$ and any two elements
of $S$ have distinct degrees. Prove that $S$ is linearly independent. There is a minor hint at the end of the assignment.
\skv

 \bf{Problem 4: }\rm 
\begin{itemize}
\item[(a)] Let $p$ be a prime, $n\geq 1$ an integer, and let $V=\dbZ_p^n$, the standard $n$-dimensional vector space
over $\dbZ_p$ (integers mod $p$). How many elements does $V$ have?
\item[(b)] Find all ordered bases for $\dbZ_2^2$. How many unordered bases does it have?
(An unordered basis is a basis in the usual sense; an ordered basis is a basis with a chosen order on its elements).
\item[(c)] Now let $p$ be any prime. How many ordered bases does $\dbZ_p^2$ have? 
\end{itemize}
{\bf Hint for (b) and (c):} Use Problem~8 from HW\#1 and one of the results stated at the beginning of this assignment.
See the end of the assignment for an additional hint for (c). 
\skv
 \bf{Problem 5: }\rm Let $F$ be a field, $X$ a set and $V=\Func(X,F)$, the set of all functions from $X$ to $F$. 
As discussed in Lecture~2, $V$ has a natural structure of a vector space over $F$. For each $x\in X$ define
the function $\delta_x\in V$ by $\delta_x(y)=0$ for all $y\neq x$ and $\delta_{x}(x)=1$. Let $S=\{\delta_x: x\in X\}$.
\begin{itemize}
\item[(a)] Prove that the set $S$ is always linearly independent.
\item[(b)] Assume that $S$ is finite. Prove that $S$ is spanning for $V$ (and hence by (a) is a basis for $V$).
\item[(c)] Now prove that if $S$ is infinite, then $S$ is not spanning for $V$. {\bf Hint:} you can find a simple
explicit function which does not lie in $\Span(V)$.
\end{itemize}
\skv
 \bf{Problem 6: }\rm Let $F$ be a field and $V=\Mat_n(F)$ the vector space of all $n\times n$
matrices over $F$. For each $1\leq i,j\leq n$ let $e_{ij}\in \Mat_n(F)$ denote the matrix
whose $(i,j)$-entry is equal to $1$ and all other entires are equal to $0$.
\begin{itemize}
\item[(a)] Prove that the set $B=\{e_{ij} : 1\leq i,j\leq n \}$ is a basis of $V$.
\item[(b)] Let $U$ be the set of upper-triangular matrices in $V$, that is,
$U$ consists of all matrices whose $(i,j)$-entry is equal to $0$ for all $i>j$.
Prove that $U$ is a subspace and find a basis of $U$ and $\dim(U)$.
{\bf Hint:} $U$ has a basis which is a subset of $B$.
\item[(c)] Now find a basis for $W=\mathfrak{sl}_n(F)$, the set of all matrices in $V$ with zero trace
(recall that $W$ was shown to be a subspace in HW\#1.2). {\bf Hint:} $W$ does not have a basis
which is a subset of $V$, but you may want to look for a basis consisting of matrices with as few nonzero entries as possible.
\end{itemize}
\skv
 \bf{Problem 7: }\rm This a continuation of Problem~1 from HW\#1. This time we are looking at $\dbQ$-subspaces of $\dbC$.
A complex number $\alpha\in\dbC$ is called {\it algebraic} if $\alpha$ is a root of a nonzero polynomial with rational coefficients,
that is, if there exist rational numbers  $q_0,\ldots, q_n$, not all $0$, such that $\sum\limits_{i=0}^n q_i \alpha^i=0$. 
A complex number is called {\it transcendental} if it is not algebraic (for example, one can show that $e$ and $\pi$ are transcendental).

Suppose you are given a specific transcendental number $\alpha\in\dbC$. Use it to construct, for every $n\in\dbN$, an explicit
$\dbQ$-subspace of $\dbC$ of dimension exactly $n$.
\newpage
{\bf Hint for 3:} Start with a concrete example,
say, $F=\dbR$, $n=2$, and $S$ consisting of $3$ polynomials of degrees $0,1$ and $2$, respectively, and look for a simple reason
why $S$ is linearly independent. Then try to generalize that argument to an arbitrary $S$ satisfying the given conditions.

\newpage
{\bf Hint for 4(c):} How many choices are there to choose the first element of an ordered basis. Once the first element has been chosen,
how may choices are there choices are there for the second element?

\end{document}
